\section{Unified backward demand analysis}
\label{sec:demand}

This section gives the finite-lattice theorem that makes DARC query-specific.
Although the construction resembles a static-analysis fixed point, the proof
is elementary and quantitative; no appeal to an infinitary fixed-point theorem
is needed.

\subsection{Finite product of resource chains}

Let \(R=\{1,\ldots,m\}\) be the finite set of resource coordinates and let
\(c_r\in\Nat\) be the capacity of coordinate \(r\).  Define
\[
  L=\prod_{r=1}^{m}\{0,1,\ldots,c_r\}.
\]
For \(d,e\in L\), write \(d\le e\) when \(d_r\le e_r\) for every \(r\), and
let \(d\join e\) be coordinatewise maximum.  Define the rank
\[
  \rank(d)=\sum_{r=1}^{m}d_r,
  \qquad C=\sum_{r=1}^{m}c_r.
\]
Thus \(0\le\rank(d)\le C\).

A transfer map \(F:L\to L\) is \emph{sound} for a plan if every requirement
reported by a live local node appears in \(F(d)\).  It is \emph{monotone} if
\(d\le e\) implies \(F(d)\le F(e)\).  Threshold-rule networks are monotone
automatically.

Given a public query seed \(q\in L\), define
\[
  d_0=q,
  \qquad
  d_{s+1}=d_s\join F(d_s).
  \tag{5.1}\label{eq:demand-iteration}
\]
This is an inflationary iteration even when \(F\) itself is not inflationary.

\begin{lemma}[Strict rank increase]
\label{lem:rank-increase}
If \(d\le e\) and \(d\ne e\), then \(\rank(d)<\rank(e)\).
\end{lemma}

\begin{proof}
For every coordinate, \(e_r-d_r\ge0\).  Since \(d\ne e\), at least one of
these differences is positive.  Hence
\[
 \rank(e)-\rank(d)=\sum_r(e_r-d_r)>0.
\]
\end{proof}

\begin{theorem}[Finite stabilization]
\label{thm:demand-stabilization}
The sequence in \eqref{eq:demand-iteration} stabilizes after at most
\(C-\rank(q)\) strict steps, and therefore after at most \(C+1\) evaluations
of the update.  Its stable value \(d^*\) satisfies \(F(d^*)\le d^*\).
\end{theorem}

\begin{proof}
The sequence is coordinatewise nondecreasing.  Whenever
\(d_{s+1}\ne d_s\), \cref{lem:rank-increase} gives
\(\rank(d_{s+1})>\rank(d_s)\).  Rank starts at \(\rank(q)\), is integral, and
is bounded by \(C\); hence there are at most \(C-\rank(q)\) strict steps.
At the first stationary stage,
\(d^*=d^*\join F(d^*)\), which is equivalent to \(F(d^*)\le d^*\).
\end{proof}

\begin{remark}
The useful bound is the sum of capacities, not the number
\(\prod_r(c_r+1)\) of states.  A plan may therefore have many Boolean resource
coordinates while still converging in a number of strict rounds linear in the
number of coordinates.
\end{remark}

\subsection{Leastness}

A demand state \(e\) is \emph{admissible above the query} if
\[
  q\le e
  \quad\text{and}\quad
  F(e)\le e.
\]
The second condition says that all requirements triggered by \(e\) are already
included in \(e\).

\begin{theorem}[Least prefixed point above the query]
\label{thm:demand-least}
Assume \(F\) is monotone.  The stable value \(d^*\) of
\eqref{eq:demand-iteration} is below every admissible state \(e\).  Therefore
it is the least prefixed point of \(F\) above \(q\).
\end{theorem}

\begin{proof}
We prove \(d_s\le e\) by induction on \(s\).  The base case is \(q\le e\).
If \(d_s\le e\), monotonicity gives \(F(d_s)\le F(e)\le e\).  Since also
\(d_s\le e\), their join satisfies
\(d_{s+1}=d_s\join F(d_s)\le e\).  At stabilization this yields
\(d^*\le e\).
\end{proof}

\begin{corollary}[Relative resource minimality]
\label{cor:relative-minimality}
No resource coordinate retained by \(d^*\) can be lowered while preserving the
query and all declared transfer summaries, unless some other coordinate is
changed so that the resulting state is no longer below \(d^*\).
\end{corollary}

The qualification ``declared transfer summaries'' is essential.  If a local
proof is replaced by a stronger lemma with a better modulus, the transfer map
changes and the new least state may be smaller.  DARC minimality is a precise
audit statement, not a claim of proof-theoretic optimality.

\subsection{Cycles are harmless when they are monotone}

Precision dependencies often contain cycles.  For instance, evaluating a
root may require an integral approximation, while the certified integration
bound may itself use a rough root bound.  Suppose
\[
  y\ge1\Rightarrow x\ge2,
  \qquad
  x\ge2\Rightarrow y\ge4,
  \qquad
  y\ge4\Rightarrow x\ge5.
\]
Starting from a query demanding \(y\ge1\), the iteration obtains the finite
closure \(x=5,y=4\).  No recursive theorem is being assumed.  The cycle occurs
in the finite resource analysis and terminates by rank.

This distinction is important:
\begin{itemize}
  \item a cycle in the demand network is solved by finite monotone closure;
  \item a cycle in mathematical execution is extracted as a loop node and
        proved by quantifier-free induction;
  \item a cycle in witness compatibility is treated as a finite relational
        problem or requires an explicit extension producer.
\end{itemize}
Conflating these three kinds of cycles produces either an unnecessary fixed
point inside arithmetic or an unjustified infinite-choice step.

\subsection{Soundness of transfer summaries}

The fixed-point theorems above are purely order-theoretic.  To connect them to
proofs, each transfer rule needs a local soundness justification.  A typical
precision rule has the form
\[
  \Prec(v)\ge p
  \quad\Longrightarrow\quad
  \Prec(u_i)\ge \rho_i(p),\;
  \Guard(g_j)=1,\;
  \Res(B_k)=1.
\]
Its soundness statement is:
\begin{quote}
Whenever the input nodes satisfy the requested precisions, the guard formulas
hold, and the listed residual producer formulas are available, the local
\(\Prf\;\PRAAx\) derivation proves the requested output contract.
\end{quote}
This is not checked merely by the metalevel transfer function.  It is tied to a
specific local theorem.  The Lean plan should package the summary and theorem
together, or generate a theorem-ledger entry linking them.

For threshold rules, the generated audit can list every firing rule at every
iteration.  The demonstration in \pathcode{examples/demo.json} records this
history.  Such a trace is useful for review but does not replace the local
object proof.

\subsection{Pruning theorem}

Let a node \(v\) have an activation coordinate \(a_v\), or let its activation
be a monotone predicate of the resource state.  Define the retained set
\[
  V^*=\{v\in V: a_v(d^*)>0\}.
\]

\begin{proposition}[Safe dead-node elimination]
\label{prop:dead-node}
Assume each node's transfer summary includes all resources and dependency
activation needed by its local proof.  Then any node outside \(V^*\) can be
removed from the proof schedule without changing the compiled conclusion or
its residual boundary.
\end{proposition}

\begin{proof}
Suppose a removed node were needed by an active node or by the public query.
By soundness of the active node's transfer summary, the removed node's
activation coordinate would be included in \(F(d^*)\).  Since
\(F(d^*)\le d^*\), that coordinate would be positive in \(d^*\), a
contradiction.  The same argument applies to a guard, residual, certificate
field, or precision needed by a retained local proof.
\end{proof}

The proposition gives a single consistency check across liveness, precision,
guards, and residuals.  It is stronger than running four independent pruning
passes and hoping they agree.

\subsection{Compositionality of demand analyses}

Suppose plans \(P\) and \(Q\) have disjoint internal coordinates and share an
explicit interface.  Their combined transfer is the join of the extended
transfers \(F_P\) and \(F_Q\).  Monotonicity is preserved.  The least closure
of the composite may be computed globally or by iterating the two component
closures to mutual stability.  This supports modular compilation: a library
lemma can publish a compact transfer summary along with its open theorem,
without exposing its internal proof graph.

One should not, however, publish a summary that forgets a hidden residual.  A
summary is part of the trusted mathematical interface in the same sense as a
modulus: an understated requirement can make the generated theorem false.
The intended Lean representation therefore couples summaries with proofs of
local contracts, while the Python reference implementation is only an audit
model.
